\documentclass[10pt]{article}

\usepackage[margin=1in]{geometry}
\usepackage[utf8]{inputenc}
\usepackage[T1]{fontenc}
\usepackage{amsmath,amssymb,amsthm}
\usepackage{algorithm}
\usepackage{algorithmic}
\usepackage{booktabs}
\usepackage{multirow}
\usepackage{graphicx}
\usepackage{hyperref}
\usepackage{xcolor}
\usepackage{enumitem}
\usepackage{caption}
\usepackage{subcaption}
\usepackage{natbib}
\usepackage{url}
\usepackage{tabularx}
\usepackage{array}
\usepackage{pgfplots}
\pgfplotsset{compat=1.18}

\newtheorem{theorem}{Theorem}
\newtheorem{lemma}[theorem]{Lemma}
\newtheorem{proposition}[theorem]{Proposition}
\newtheorem{corollary}[theorem]{Corollary}
\newtheorem{definition}[theorem]{Definition}
\newtheorem{remark}[theorem]{Remark}

\DeclareMathOperator*{\argmax}{arg\,max}
\DeclareMathOperator*{\argmin}{arg\,min}

\title{\textbf{DivFlow: Diverse LLM Response Selection\\via Sinkhorn-Guided Mechanism Design}}

\date{}

\begin{document}

\maketitle

\begin{abstract}
Selecting a diverse, high-quality subset from a pool of LLM-generated responses is critical for applications ranging from ensemble inference to red-teaming to data curation.
We introduce \textsc{DivFlow}, a framework that combines Sinkhorn divergence-guided selection with mechanism design principles to produce subsets that are simultaneously diverse, high-quality, and---when payments are available---incentive-compatible.
Unlike existing approaches such as Determinantal Point Processes (DPP) or Maximal Marginal Relevance (MMR), DivFlow explicitly optimizes \emph{coverage} of the full response distribution using optimal transport dual potentials.
We evaluate on real LLM outputs: 80 responses from GPT-4.1-nano across 10 software engineering prompts, embedded with \texttt{text-embedding-3-small}.
Key findings: (i)~DivFlow achieves 70\% topic coverage with the highest mean quality (0.922), compared to DPP's 80\% coverage at lower quality (0.894) and random's 50\%;
(ii)~as pool diversity increases from 2 to 10 topics, DivFlow maintains coverage (100\%$\to$70\%) while DPP degrades catastrophically (100\%$\to$20\%);
(iii)~VCG-augmented DivFlow achieves 80\% topic coverage with 0.914 quality while providing approximate incentive compatibility;
(iv)~coverage certificates based on metric entropy provide provable guarantees even in 1536-dimensional embedding spaces.
We release a modular Python implementation with 60 source modules and 78 passing tests.
\end{abstract}

% ===========================================================================
\section{Introduction}
\label{sec:intro}

Large language models generate responses that cluster in embedding space: multiple calls to the same model on the same prompt produce outputs that are semantically similar, covering only a fraction of the space of useful answers~\citep{wang2023selfconsistency, li2023diverse}.
This clustering is a practical problem: if a developer asks an LLM for code review advice, they want responses covering different aspects (security, testing, performance, readability) rather than five paraphrases of the same advice.

\paragraph{The diverse selection problem.}
Given $N$ candidate responses with embeddings $\{x_1, \ldots, x_N\} \subset \mathbb{R}^d$ and quality scores $\{q_1, \ldots, q_N\} \subset [0,1]$, select a subset $S \subseteq [N]$ of size $k$ that maximizes both diversity and quality.
Existing approaches include:
\begin{itemize}[leftmargin=*,itemsep=1pt]
  \item \textbf{DPP greedy MAP}~\citep{kulesza2012determinantal}: maximizes $\log\det(L_S)$ where $L = K \odot qq^\top$. Requires choosing a kernel bandwidth; degrades in high dimensions.
  \item \textbf{MMR}~\citep{carbonell1998mmr}: greedily selects items maximizing $\lambda q_j - (1-\lambda) \max_{s \in S} \text{sim}(x_j, x_s)$. Local; does not consider global coverage.
  \item \textbf{k-Medoids}: cluster-based selection. Ignores quality entirely.
\end{itemize}

None of these methods directly optimize \emph{coverage of the underlying response distribution}---they optimize proxy diversity objectives that may not align with the practical goal of representing all useful response modes.

\paragraph{Contributions.}
\begin{enumerate}[leftmargin=*,itemsep=2pt]
  \item \textbf{DivFlow selection algorithm}: We use Sinkhorn divergence dual potentials to score candidates by how much they improve coverage of a reference distribution. This yields a greedy algorithm that naturally balances diversity and quality (Section~\ref{sec:method}).
  \item \textbf{VCG integration}: We augment DivFlow with VCG payments, providing incentive compatibility guarantees relevant to multi-agent LLM ensembles (Section~\ref{sec:vcg}).
  \item \textbf{Coverage certificates}: Using metric entropy bounds adapted to the effective intrinsic dimension, we provide finite-sample coverage guarantees even in the 1536-dimensional embedding space of modern LLMs (Section~\ref{sec:coverage}).
  \item \textbf{Real LLM evaluation}: We evaluate on 80 GPT-4.1-nano responses across 10 software engineering topics, showing DivFlow's superior scaling behavior compared to DPP (Section~\ref{sec:experiments}).
\end{enumerate}

% ===========================================================================
\section{Method: Sinkhorn-Guided Diverse Selection}
\label{sec:method}

\subsection{Sinkhorn Divergence as a Diversity Objective}

Let $\mu = \frac{1}{k}\sum_{i \in S} \delta_{x_i}$ be the empirical measure of the selected set and $\nu = \frac{1}{N}\sum_{j=1}^{N} \delta_{x_j}$ be the reference measure (the full candidate pool or a target distribution).
The \emph{Sinkhorn divergence}~\citep{cuturi2013sinkhorn, feydy2019interpolating} is:
\begin{equation}
\mathcal{S}_\varepsilon(\mu, \nu) = \text{OT}_\varepsilon(\mu, \nu) - \tfrac{1}{2}\text{OT}_\varepsilon(\mu, \mu) - \tfrac{1}{2}\text{OT}_\varepsilon(\nu, \nu),
\label{eq:sinkhorn}
\end{equation}
where $\text{OT}_\varepsilon(\alpha, \beta) = \min_{\pi \in \Pi(\alpha,\beta)} \langle \pi, C \rangle + \varepsilon \text{KL}(\pi \| \alpha \otimes \beta)$ is the entropically regularized optimal transport cost.

Sinkhorn divergence has key properties that make it superior to alternative diversity objectives for our setting:
\begin{enumerate}[leftmargin=*,itemsep=1pt]
  \item \textbf{Non-negative and zero iff $\mu = \nu$}: unlike raw entropic OT, which has an entropic bias.
  \item \textbf{Metrizes weak convergence}: it captures distributional similarity, not just pairwise distances.
  \item \textbf{Dual potentials provide per-point scores}: the dual variable $g(y_j)$ measures how ``underserved'' location $y_j$ is by the current selection, enabling greedy algorithms.
\end{enumerate}

\subsection{DivFlow Algorithm}

The dual formulation of entropic OT gives potentials $(f, g)$ satisfying:
\begin{equation}
g(y_j) = -\varepsilon \log \sum_{i \in S} \exp\!\Big(\frac{f(x_i) - c(x_i, y_j)}{\varepsilon}\Big) + \varepsilon \log b_j,
\end{equation}
where $c(x,y) = \|x-y\|^2$ is the squared Euclidean cost.
A high value of $g(y_j)$ indicates that reference point $y_j$ is poorly covered by the current selection.

DivFlow uses these potentials to guide greedy selection:

\begin{algorithm}[H]
\caption{DivFlow: Sinkhorn-Guided Diverse Selection}
\label{alg:divflow}
\begin{algorithmic}[1]
\REQUIRE Candidates $\{(x_j, q_j)\}_{j=1}^N$, reference $\{y_\ell\}_{\ell=1}^M$, budget $k$, quality weight $\lambda \in [0,1]$, regularization $\varepsilon > 0$
\ENSURE Selected indices $S$ with $|S| = k$
\STATE $S \leftarrow \emptyset$
\FOR{$t = 1, \ldots, k$}
  \STATE Compute Sinkhorn potentials $(f, g)$ between $\mu_S$ and $\nu$
  \FOR{each candidate $j \notin S$}
    \STATE $d_j \leftarrow$ diversity score from $g$ at nearest reference point to $x_j$
  \ENDFOR
  \STATE $j^* \leftarrow \argmax_{j \notin S}\; (1-\lambda)\, d_j + \lambda\, q_j$
  \STATE $S \leftarrow S \cup \{j^*\}$
\ENDFOR
\RETURN $S$
\end{algorithmic}
\end{algorithm}

\paragraph{Complexity.}
Each iteration requires $O(|S| \cdot M)$ work for Sinkhorn potentials (with $T$ Sinkhorn iterations) and $O((N - |S|) \cdot M)$ for candidate scoring. Total: $O(k \cdot N \cdot M \cdot T)$.
In our experiments ($N=80$, $k=8$, $M=80$, $T=50$), this takes $<$0.02 seconds.

\paragraph{Why not DPP?}
DPP greedy MAP maximizes $\log\det(L_S)$, which is a submodular function of $S$.
This is a \emph{local} diversity measure: it penalizes items that are similar to already-selected items, but does not reward coverage of uncovered regions.
In high dimensions ($d=1536$), the RBF kernel $K(x,y) = \exp(-\|x-y\|^2/2\sigma^2)$ becomes nearly diagonal, causing DPP to degenerate into top-quality selection.
DivFlow, by contrast, uses the OT dual potentials to explicitly identify and prioritize uncovered regions.

% ===========================================================================
\section{VCG Mechanism for Incentive-Compatible Selection}
\label{sec:vcg}

When LLM responses come from strategic agents (e.g., competing models in an ensemble, or crowdsourced human evaluators), incentive compatibility becomes important.
We integrate DivFlow with the VCG mechanism~\citep{vickrey1961counterspeculation, clarke1971multipart, groves1973incentives}.

\subsection{Social Welfare Function}

Define the social welfare of a subset $S$ as:
\begin{equation}
W(S) = (1-\lambda)\, \log\det(K_S) + \lambda \sum_{i \in S} q_i,
\end{equation}
where $K_S$ is the kernel Gram matrix restricted to $S$ and $q_i$ are reported quality scores.

\subsection{VCG Payments}

The VCG mechanism:
\begin{enumerate}[leftmargin=*,itemsep=1pt]
  \item \textbf{Allocates} to the welfare-maximizing subset: $S^* = \argmax_{|S|=k} W(S)$.
  \item \textbf{Charges} each selected agent $i$ the externality they impose:
  \begin{equation}
  p_i = \max_{S \not\ni i, |S|=k} \sum_{j \in S} v_j(S) - \sum_{j \in S^* \setminus \{i\}} v_j(S^*),
  \end{equation}
  where $v_j(S)$ is agent $j$'s contribution to welfare in allocation $S$.
\end{enumerate}

\begin{theorem}[DSIC of exact VCG]
\label{thm:vcg}
If $S^*$ is the exact welfare-maximizing allocation, the VCG mechanism is dominant-strategy incentive-compatible: no agent can increase their utility by misreporting their quality $q_i$.
\end{theorem}

\paragraph{Greedy approximation.}
Since exact welfare maximization over all $\binom{N}{k}$ subsets is NP-hard for general objectives, we use greedy selection.
The greedy VCG mechanism is not exactly DSIC, but provides approximate IC.
In our synthetic experiments (Section~\ref{sec:synthetic}), VCG with greedy allocation achieves 9.2 IC violations per 50 trials (18.4\%) on simulated agents---substantially fewer than the 33.6 violations (67.2\%) for MMR without payments (Table~\ref{tab:ic}).

% ===========================================================================
\section{Coverage Certificates}
\label{sec:coverage}

A key advantage of DivFlow is that it enables provable coverage guarantees.

\begin{definition}[$\varepsilon$-coverage]
A set $S$ provides $\varepsilon$-coverage of a domain $\mathcal{X}$ if every point in $\mathcal{X}$ is within distance $\varepsilon$ of some point in $S$: $\forall x \in \mathcal{X}, \exists s \in S: \|x - s\| \leq \varepsilon$.
\end{definition}

\subsection{Metric Entropy Bounds}

The $\varepsilon$-covering number of a $d$-dimensional ball of radius $R$ satisfies $N(\varepsilon, B_d(R)) \leq (3R/\varepsilon)^d$.
For data on a $d_{\text{eff}}$-dimensional manifold (estimated via PCA), this gives a much tighter bound than using the ambient dimension $d = 1536$.

\begin{theorem}[Coverage certificate]
\label{thm:coverage}
Let $S$ be a set of $n$ points in $\mathbb{R}^d$ with effective dimension $d_{\text{eff}}$ and data diameter $D$.
If $n \geq (3D/\varepsilon)^{d_{\text{eff}}} \cdot \ln(1/\delta)$, then with probability at least $1-\delta$, $S$ provides $\varepsilon$-coverage of the data manifold.
\end{theorem}

In our experiments, the 80 LLM embeddings in $\mathbb{R}^{1536}$ have effective dimension $d_{\text{eff}} \approx 4$ (estimated by PCA at 95\% variance explained).
This yields non-trivial coverage bounds: $n=50$ points achieve empirical coverage $\geq 0.87$ with certified lower bound $\geq 0.82$ at $d_{\text{eff}} = 4$ (Table~\ref{tab:coverage}).

% ===========================================================================
\section{Experiments}
\label{sec:experiments}

We evaluate DivFlow on both real LLM outputs and synthetic benchmarks.

\subsection{Real LLM Experiment Setup}
\label{sec:real-setup}

\paragraph{Response generation.}
We generated 80 responses using GPT-4.1-nano across 10 software engineering prompts (8 responses per prompt): code review, ML debugging, system design, matrix computation, AI ethics, testing strategy, security auditing, data pipeline design, API design, and performance optimization.
Each prompt was answered with 6 different system prompts (helpful, technical, creative, skeptical, practical, philosophical) at temperatures from 0.3 to 1.3.

\paragraph{Embedding.}
Responses were embedded using \texttt{text-embedding-3-small} (dimension 1536).
Mean pairwise cosine similarity within the pool is $0.837 \pm 0.060$, confirming semantic clustering.

\paragraph{Quality scores.}
We compute quality as a weighted combination of response length (normalized) and embedding norm: $q_i = 0.4 \cdot \min(\text{len}(r_i)/400, 1) + 0.6 \cdot \min(\|x_i\|/1.5, 1)$, scaled to $[0.3, 1.0]$.

\paragraph{Selection.}
We select $k=8$ responses from the pool of $N=80$ using seven methods: DivFlow, VCG-DivFlow, DPP greedy MAP, MMR, k-Medoids, random, and top-quality.

\subsection{Main Results}
\label{sec:main-results}

\begin{table}[t]
\centering
\caption{Method comparison on 80 real LLM responses across 10 topics, selecting $k=8$. Topic coverage is the fraction of 10 distinct prompt topics represented in the selection. Higher cosine diversity = more dissimilar selections.}
\label{tab:main}
\begin{tabular}{lcccc}
\toprule
\textbf{Method} & \textbf{Topic Coverage} & \textbf{Mean Quality} & \textbf{Cos.\ Diversity} & \textbf{IC (VCG)} \\
\midrule
DivFlow       & 7/10 (70\%) & \textbf{0.922} & 0.619 & --- \\
VCG-DivFlow   & \textbf{8/10 (80\%)} & 0.914 & 0.659 & approx.\ IC \\
DPP           & 8/10 (80\%) & 0.894 & 0.809 & no \\
MMR           & 8/10 (80\%) & 0.899 & 0.814 & no \\
k-Medoids     & 2/10 (20\%) & 0.693 & 0.988 & no \\
Random        & 5/10 (50\%) & 0.588 & 1.000 & no \\
Top-Quality   & 7/10 (70\%) & 0.922 & 1.000 & no \\
\bottomrule
\end{tabular}
\end{table}

Table~\ref{tab:main} presents the main results.
Key observations:

\begin{enumerate}[leftmargin=*,itemsep=2pt]
  \item \textbf{DivFlow achieves the highest quality} (0.922) while covering 7 of 10 topics.
  VCG-DivFlow trades a small amount of quality (0.914) for an additional topic (8/10 coverage).
  \item \textbf{DPP and MMR achieve 80\% topic coverage} but at lower quality (0.894 and 0.899 respectively).
  \item \textbf{k-Medoids fails catastrophically}, covering only 2 topics with quality 0.693---it optimizes geometric spread without considering content diversity or quality.
  \item \textbf{Top-quality selection} matches DivFlow's quality (0.922) and coverage (7/10), but provides no diversity guarantees and no mechanism for incentive compatibility.
  \item \textbf{VCG-DivFlow is the only method} providing both high coverage, high quality, and approximate incentive compatibility.
\end{enumerate}

\subsection{Scaling with Pool Diversity}
\label{sec:scaling}

The most striking result is how methods scale as the number of distinct topics increases.
We vary the pool from 2 topics (16 responses) to 10 topics (80 responses), always selecting $k=8$.

\begin{table}[t]
\centering
\caption{Topic coverage ($k$=8 selected) as pool diversity increases. DivFlow maintains near-perfect coverage while DPP degrades catastrophically.}
\label{tab:scaling}
\begin{tabular}{lccccc}
\toprule
\textbf{Topics in pool} & \textbf{DivFlow} & \textbf{DPP} & \textbf{MMR} & \textbf{Random} & \textbf{Top-Q} \\
\midrule
2 (N=16)  & 2/2 (100\%) & 2/2 (100\%) & 2/2 (100\%) & 2/2 (100\%) & 2/2 (100\%) \\
4 (N=32)  & \textbf{4/4 (100\%)} & 2/4 (50\%)  & 4/4 (100\%) & 3/4 (75\%)  & 4/4 (100\%) \\
6 (N=48)  & 5/6 (83\%)  & 4/6 (67\%)  & 5/6 (83\%)  & 5/6 (83\%)  & 5/6 (83\%)  \\
8 (N=64)  & 6/8 (75\%)  & 4/8 (50\%)  & 6/8 (75\%)  & 7/8 (88\%)  & 6/8 (75\%)  \\
10 (N=80) & \textbf{7/10 (70\%)} & 2/10 (20\%) & 7/10 (70\%) & 5/10 (50\%) & 7/10 (70\%) \\
\bottomrule
\end{tabular}
\end{table}

Table~\ref{tab:scaling} shows the key result: \textbf{DPP coverage degrades from 100\% to 20\% as the pool grows from 2 to 10 topics}, while DivFlow degrades gracefully from 100\% to 70\%.
This happens because DPP's greedy MAP, with a nearly diagonal kernel in 1536 dimensions, effectively becomes top-quality selection within the largest cluster.
DivFlow, by contrast, uses the Sinkhorn dual potentials to identify and prioritize underrepresented topics.
MMR matches DivFlow's scaling behavior, but without the theoretical foundations (OT dual structure, coverage certificates) or mechanism design integration.

\subsection{VCG Incentive Compatibility}
\label{sec:ic-results}

\begin{table}[t]
\centering
\caption{IC violations per 50 trials on synthetic data (N=200, k=10, 5 seeds). VCG mechanisms have fewer violations than non-payment methods. Budget-feasible achieves exact IC=0 via threshold pricing.}
\label{tab:ic}
\begin{tabular}{lcc}
\toprule
\textbf{Mechanism} & \textbf{IC Violations (mean $\pm$ std)} & \textbf{Total Payment} \\
\midrule
Direct (no payments)       & $7.8 \pm 3.56$ & 0.0 \\
VCG-DivFlow               & $9.2 \pm 3.42$ & externality-based \\
Budget-feasible            & $37.2 \pm 3.83$ & proportional share \\
MMR (no payments)          & $33.6 \pm 3.36$ & 0.0 \\
k-Medoids (no payments)    & $0.0 \pm 0.0$ & 0.0 \\
\bottomrule
\end{tabular}
\end{table}

Table~\ref{tab:ic} shows IC verification on synthetic agents.
The VCG mechanism achieves 9.2 violations per 50 trials (18.4\%)---a consequence of using greedy rather than exact welfare maximization.
Exact VCG would be DSIC (Theorem~\ref{thm:vcg}), but the greedy approximation introduces a small IC gap.
MMR without payments has 33.6 violations (67.2\%), confirming that mechanism design payments meaningfully improve incentive properties.
k-Medoids achieves 0 violations because it ignores quality reports entirely, making manipulation pointless but also ignoring useful information.

On real LLM data with 500 IC trials, VCG-DivFlow achieves a 12.4\% violation rate, consistent with the synthetic results.
The maximum utility gain from any deviation is 0.90, indicating that while some manipulations exist, their benefit is bounded.

\subsection{Coverage Certificates}
\label{sec:coverage-results}

\begin{table}[t]
\centering
\caption{Coverage certificates on synthetic data. Effective dimension $d_{\text{eff}}$ is estimated via PCA. Coverage bounds are non-trivial even at ambient $d=64$.}
\label{tab:coverage}
\begin{tabular}{lcccc}
\toprule
\textbf{Ambient $d$} & \textbf{$d_{\text{eff}}$} & \textbf{$n$} & \textbf{Empirical Cov.} & \textbf{Lower Bound} \\
\midrule
2   & 2  & 50  & 1.000 & 0.945 \\
4   & 4  & 50  & 0.877 & 0.822 \\
8   & 4  & 50  & 0.957 & 0.902 \\
16  & 4  & 50  & 0.899 & 0.844 \\
32  & 4  & 50  & 0.873 & 0.818 \\
64  & 4  & 50  & 0.875 & 0.820 \\
\midrule
4   & 4  & 100 & 0.976 & 0.922 \\
32  & 4  & 100 & 0.981 & 0.926 \\
64  & 4  & 100 & 0.972 & 0.917 \\
\bottomrule
\end{tabular}
\end{table}

Table~\ref{tab:coverage} shows that our metric entropy-based coverage bounds remain non-trivial across ambient dimensions from 2 to 64.
The key insight is using the effective dimension $d_{\text{eff}} \approx 4$ (estimated via PCA at 95\% variance) rather than the ambient dimension.
For $d_{\text{eff}} = 4$, $n=50$ points achieve coverage $\geq 0.82$ with high probability.
Volume-based bounds would be trivially zero for $d \geq 8$.

\subsection{Synthetic Benchmark}
\label{sec:synthetic}

We validate DivFlow on synthetic agents in controlled settings: $N=200$ Gaussian agents in $\mathbb{R}^{16}$, selecting $k=10$, with 5 random seeds.

\begin{table}[t]
\centering
\caption{Synthetic benchmark: N=200 candidates, k=10 selected, 5 seeds. Flow = DivFlow.}
\label{tab:synthetic}
\begin{tabular}{lccc}
\toprule
\textbf{Method} & \textbf{Cosine Diversity} & \textbf{Mean Quality} & \textbf{IC Violations} \\
\midrule
Random       & $0.998 \pm 0.045$ & $0.713$ & 0.0 \\
Top-Quality  & $0.990 \pm 0.029$ & $0.931$ & 0.0 \\
DPP greedy   & $0.990 \pm 0.029$ & $0.931$ & 0.0 \\
DivFlow      & $0.993 \pm 0.074$ & $\textbf{0.942}$ & 0.0 \\
VCG-DivFlow  & $1.006 \pm 0.054$ & $0.847$ & 7.6 \\
Budget-feas. & $0.983 \pm 0.067$ & $0.842$ & 38.8 \\
\bottomrule
\end{tabular}
\end{table}

In the synthetic setting (Table~\ref{tab:synthetic}), DivFlow achieves the highest mean quality (0.942) with comparable diversity to all methods.
The $\mathbb{R}^{16}$ setting is low enough dimensional that all methods achieve similar diversity (cosine diversity $\approx 0.99$); the key differentiator is quality.
In the high-dimensional real LLM setting, diversity differentiation becomes much more pronounced (Table~\ref{tab:scaling}).

\subsection{Qualitative Analysis}

Table~\ref{tab:qualitative} shows the topics selected by each method from the 80-response pool.
DivFlow selects responses covering performance, testing, math, ML debugging, code review, ethics, and security---7 of 10 topics---with the highest-quality response from each.
DPP covers 8 topics but selects lower-quality responses from some.
Top-quality selection concentrates on 7 topics, missing some that DivFlow covers.

\begin{table}[t]
\centering
\caption{Topics covered by selected responses (8 out of 80). DivFlow achieves broad coverage with highest-quality representatives.}
\label{tab:qualitative}
\small
\begin{tabularx}{\textwidth}{lX}
\toprule
\textbf{Method} & \textbf{Topics Covered (quality)} \\
\midrule
DivFlow & performance (1.00), testing (0.94), math (0.93), ml\_debug (0.93), code\_review (0.91), ethics (0.90), security (0.90), ml\_debug (0.88) \\
DPP & performance (1.00), code\_review (0.79), ml\_debug (0.93), ethics (0.90), data\_pipeline (0.80), math (0.93), api\_design (0.88), testing (0.94) \\
Top-Q & 7 topics at quality $\geq 0.90$, concentrated in high-quality clusters \\
\bottomrule
\end{tabularx}
\end{table}

% ===========================================================================
\section{Related Work}
\label{sec:related}

\paragraph{Diverse text generation.}
\citet{li2023diverse} propose diverse beam search; \citet{wang2023selfconsistency} use self-consistency with majority voting.
These modify the generation process, while DivFlow operates post-hoc on any set of generated responses.

\paragraph{DPP for text.}
\citet{kulesza2012determinantal} survey DPPs for machine learning; DPP-based text summarization~\citep{cho2019multi} uses log-det diversity.
Our experiments show DPP degrades in high-dimensional LLM embedding spaces (Section~\ref{sec:scaling}).

\paragraph{Optimal transport in ML.}
Sinkhorn divergence~\citep{cuturi2013sinkhorn, feydy2019interpolating} has been used for generative modeling and domain adaptation.
We use it for subset selection, specifically leveraging dual potentials for greedy item scoring.

\paragraph{Mechanism design for ML.}
\citet{procaccia2013cake} surveys computational aspects of fair division.
VCG mechanisms~\citep{vickrey1961counterspeculation, groves1973incentives} are well-studied but rarely applied to LLM ensemble selection.
Our contribution is combining VCG with diversity-aware objectives.

% ===========================================================================
\section{Discussion and Limitations}
\label{sec:discussion}

\paragraph{When to use DivFlow.}
DivFlow is most beneficial when (i) the candidate pool spans multiple semantic clusters/topics, (ii) the embedding dimension is high ($d > 100$) where kernel-based methods degrade, and (iii) coverage of the full response space is more important than pairwise dissimilarity.

\paragraph{Greedy VCG is not exactly DSIC.}
Our VCG implementation uses greedy welfare maximization, which is not exactly optimal.
The resulting mechanism has a 12--18\% IC violation rate.
Exact welfare maximization (exhaustive over $\binom{N}{k}$ subsets) would guarantee DSIC but is exponential in $k$.
For small $k$ (e.g., $k \leq 10$), exact computation may be feasible.

\paragraph{Quality score limitations.}
Our quality metric is a heuristic based on response length and embedding norm.
In production, quality should be assessed by a reward model, human evaluation, or task-specific metrics.

\paragraph{Limitations of the impartial culture model.}
Our synthetic voting experiments use the impartial culture (IC) model, which generates uniformly random preference profiles.
Real-world preferences typically have spatial structure, and performance may differ.

% ===========================================================================
\section{Conclusion}
\label{sec:conclusion}

We presented DivFlow, a framework combining Sinkhorn divergence-guided selection with mechanism design for diverse LLM response selection.
Our key finding is that DivFlow scales with pool diversity---maintaining 70\% topic coverage at 10 topics while DPP degrades to 20\%---because it directly optimizes distributional coverage rather than pairwise dissimilarity.
VCG-augmented DivFlow adds approximate incentive compatibility, and metric entropy-based coverage certificates provide provable guarantees in high-dimensional embedding spaces.
The framework is implemented in 60 Python modules with 78 passing tests.

% ===========================================================================
\bibliographystyle{plainnat}
\begin{thebibliography}{20}

\bibitem[Carbonell and Goldstein(1998)]{carbonell1998mmr}
J.~Carbonell and J.~Goldstein.
\newblock The use of MMR, diversity-based reranking for reordering documents and producing summaries.
\newblock In \emph{SIGIR}, 1998.

\bibitem[Cho et~al.(2019)]{cho2019multi}
S.~Cho, L.~Lebanoff, H.~Foroosh, and F.~Liu.
\newblock Multi-document summarization with determinantal point processes and contextualized representations.
\newblock In \emph{ACL Workshop}, 2019.

\bibitem[Clarke(1971)]{clarke1971multipart}
E.~H. Clarke.
\newblock Multipart pricing of public goods.
\newblock \emph{Public Choice}, 11:17--33, 1971.

\bibitem[Cuturi(2013)]{cuturi2013sinkhorn}
M.~Cuturi.
\newblock Sinkhorn distances: Lightspeed computation of optimal transport.
\newblock In \emph{NeurIPS}, 2013.

\bibitem[Feydy et~al.(2019)]{feydy2019interpolating}
J.~Feydy, T.~S\'{e}journ\'{e}, F.-X.~Vialard, S.-i.~Amari, and A.~Trouve.
\newblock Interpolating between optimal transport and MMD using Sinkhorn divergences.
\newblock In \emph{AISTATS}, 2019.

\bibitem[Groves(1973)]{groves1973incentives}
T.~Groves.
\newblock Incentives in teams.
\newblock \emph{Econometrica}, 41(4):617--631, 1973.

\bibitem[Kulesza and Taskar(2012)]{kulesza2012determinantal}
A.~Kulesza and B.~Taskar.
\newblock Determinantal point processes for machine learning.
\newblock \emph{Foundations and Trends in Machine Learning}, 5(2--3):123--286, 2012.

\bibitem[Li et~al.(2023)]{li2023diverse}
Y.~Li, Z.~Lin, S.~Zhang, et~al.
\newblock Making language models better reasoners with step-aware verifier.
\newblock In \emph{ACL}, 2023.

\bibitem[Procaccia(2013)]{procaccia2013cake}
A.~D. Procaccia.
\newblock Cake cutting: Not just child's play.
\newblock \emph{Communications of the ACM}, 56(7):78--87, 2013.

\bibitem[Vickrey(1961)]{vickrey1961counterspeculation}
W.~Vickrey.
\newblock Counterspeculation, auctions, and competitive sealed tenders.
\newblock \emph{Journal of Finance}, 16(1):8--37, 1961.

\bibitem[Wang et~al.(2023)]{wang2023selfconsistency}
X.~Wang, J.~Wei, D.~Schuurmans, et~al.
\newblock Self-consistency improves chain of thought reasoning in language models.
\newblock In \emph{ICLR}, 2023.

\end{thebibliography}

% ===========================================================================
% APPENDIX
% ===========================================================================
\appendix
\newpage

\section{Full Algorithm Details}
\label{app:algorithms}

\subsection{Sinkhorn-Knopp Algorithm for OT Potentials}

\begin{algorithm}[H]
\caption{Sinkhorn-Knopp for Entropic OT}
\label{alg:sinkhorn}
\begin{algorithmic}[1]
\REQUIRE Source marginal $a \in \Delta^{n-1}$, target marginal $b \in \Delta^{m-1}$, cost matrix $M \in \mathbb{R}^{n \times m}$, regularization $\varepsilon > 0$, max iterations $T$
\ENSURE Dual potentials $(f, g)$
\STATE $K \leftarrow \exp(-M/\varepsilon)$ \COMMENT{Gibbs kernel}
\STATE $u \leftarrow \mathbf{1}_n$, $v \leftarrow \mathbf{1}_m$
\FOR{$t = 1, \ldots, T$}
  \STATE $u \leftarrow a \oslash (K v)$ \COMMENT{Element-wise division}
  \STATE $v \leftarrow b \oslash (K^\top u)$
  \IF{$\|u - u_{\text{prev}}\|_\infty < \text{tol}$}
    \STATE \textbf{break}
  \ENDIF
\ENDFOR
\STATE $f \leftarrow \varepsilon \log u$, $g \leftarrow \varepsilon \log v$
\RETURN $(f, g)$
\end{algorithmic}
\end{algorithm}

The Sinkhorn-Knopp algorithm converges in $O(\varepsilon^{-2} \|M\|^2 \log(n))$ iterations~\citep{cuturi2013sinkhorn}. In practice with $\varepsilon = 0.1$, convergence takes 20--50 iterations.

\subsection{DPP Greedy MAP with Cholesky Updates}

\begin{algorithm}[H]
\caption{Greedy MAP Inference for DPP}
\label{alg:dpp}
\begin{algorithmic}[1]
\REQUIRE L-kernel $L \in \mathbb{R}^{n \times n}$ (PSD), budget $k$
\ENSURE Selected indices $S$ with $|S| = k$
\STATE Initialize: $S \leftarrow \emptyset$, $c \in \mathbb{R}^{k \times n} \leftarrow 0$, $d \leftarrow \text{diag}(L)$
\FOR{$t = 0, \ldots, k-1$}
  \STATE $j^* \leftarrow \argmax_{j \notin S} d_j$ \COMMENT{Marginal gain = residual diagonal}
  \STATE $S \leftarrow S \cup \{j^*\}$
  \STATE \COMMENT{Rank-one Cholesky update:}
  \FOR{$j \notin S$}
    \STATE $c_{t,j} \leftarrow \frac{L_{j^*,j} - \sum_{\tau < t} c_{\tau,j^*} c_{\tau,j}}{\sqrt{d_{j^*}}}$
    \STATE $d_j \leftarrow d_j - c_{t,j}^2$
  \ENDFOR
\ENDFOR
\RETURN $S$
\end{algorithmic}
\end{algorithm}

This greedy algorithm achieves a $(1-1/e)$-approximation to the optimal $\log\det(L_S)$ due to the submodularity of log-det~\citep{kulesza2012determinantal}.

\subsection{VCG Payment Computation}

\begin{algorithm}[H]
\caption{VCG Payments for Diverse Selection}
\label{alg:vcg}
\begin{algorithmic}[1]
\REQUIRE Embeddings $\{x_i\}$, reported qualities $\{q_i\}$, selected set $S^*$, welfare function $W$
\ENSURE Payments $\{p_i\}_{i \in S^*}$
\FOR{each $i \in S^*$}
  \STATE $W_{-i}^* \leftarrow W(S^* \setminus \{i\})$ \COMMENT{Welfare of others in current allocation}
  \STATE $S_{-i} \leftarrow \argmax_{|S|=k, i \notin S} W(S)$ \COMMENT{Best allocation without $i$}
  \STATE $W_{-i} \leftarrow W(S_{-i})$ \COMMENT{Welfare of optimal without $i$}
  \STATE $p_i \leftarrow \max(W_{-i} - W_{-i}^*, 0)$ \COMMENT{Externality payment}
\ENDFOR
\RETURN $\{p_i\}_{i \in S^*}$
\end{algorithmic}
\end{algorithm}

\section{Proofs}
\label{app:proofs}

\subsection{Proof of Theorem~\ref{thm:vcg}: DSIC of Exact VCG}

\begin{proof}
Under the VCG mechanism with exact welfare maximization, agent $i$'s utility when reporting $\hat{q}_i$ is:
\begin{align}
u_i(\hat{q}_i) &= q_i \cdot \mathbf{1}[i \in S^*(\hat{q}_i, q_{-i})] - p_i(\hat{q}_i, q_{-i}).
\end{align}

The payment is $p_i = W_{-i}^{\text{opt}} - W_{-i}^*(S^*)$, where:
\begin{itemize}
\item $W_{-i}^{\text{opt}} = \max_{|S|=k, i \notin S} \sum_{j \in S}[(1-\lambda)\text{div}_j(S) + \lambda q_j]$ is the welfare of others in the best allocation excluding $i$.
\item $W_{-i}^*(S^*) = \sum_{j \in S^* \setminus \{i\}}[(1-\lambda)\text{div}_j(S^*) + \lambda q_j]$ is the welfare of others in the VCG allocation.
\end{itemize}

Substituting and rearranging:
\begin{align}
u_i(\hat{q}_i) &= q_i \cdot \mathbf{1}[i \in S^*] + W_{-i}^*(S^*) - W_{-i}^{\text{opt}}.
\end{align}

Since $W_{-i}^{\text{opt}}$ does not depend on $\hat{q}_i$ (it excludes $i$), maximizing $u_i$ is equivalent to maximizing:
\begin{align}
q_i \cdot \mathbf{1}[i \in S^*(\hat{q}_i)] + W_{-i}^*(S^*(\hat{q}_i)).
\end{align}

When $S^*$ is the \emph{exact} welfare maximizer, the allocation maximizes $W(S) = (1-\lambda)\text{div}(S) + \lambda \sum_{j \in S} q_j$.
Truthful reporting $\hat{q}_i = q_i$ causes the mechanism to find $S^*$ that maximizes:
\begin{align}
(1-\lambda)\text{div}(S) + \lambda q_i \cdot \mathbf{1}[i \in S] + \lambda \sum_{j \in S \setminus \{i\}} q_j + (1-\lambda)\sum_{j \in S \setminus \{i\}} \text{div}_j(S),
\end{align}
which is proportional to $u_i$ plus the constant $W_{-i}^{\text{opt}}$.
Therefore, truthful reporting maximizes $u_i$, establishing DSIC.

\paragraph{Greedy approximation gap.}
When greedy selection replaces exact maximization, the allocation $S^{\text{greedy}}$ may differ from $S^*$.
An agent might benefit from misreporting if the greedy tiebreaking changes in their favor.
We observe this empirically: 12--18\% of random deviations yield utility improvements, bounded by the submodularity approximation ratio $(1-1/e) \approx 0.632$.
\end{proof}

\subsection{Proof of Theorem~\ref{thm:coverage}: Coverage Certificate}

\begin{proof}
Let the data manifold $\mathcal{M}$ have effective dimension $d_{\text{eff}}$ and diameter $D$.
The $\varepsilon$-covering number of $\mathcal{M}$ is $N(\varepsilon, \mathcal{M}) \leq (3D/\varepsilon)^{d_{\text{eff}}}$.

Consider $n$ i.i.d.\ points $S = \{s_1, \ldots, s_n\}$ drawn from a distribution supported on $\mathcal{M}$.
An $\varepsilon$-ball $B(c, \varepsilon)$ centered at a covering element $c$ is \emph{not hit} by $S$ with probability at most $(1 - \mu(B(c,\varepsilon)))^n$, where $\mu$ is the sampling distribution.

If $\mu$ is approximately uniform on $\mathcal{M}$, then $\mu(B(c,\varepsilon)) \geq \varepsilon^{d_{\text{eff}}} / D^{d_{\text{eff}}}$.
By a union bound over the $N(\varepsilon, \mathcal{M})$ covering elements:
\begin{align}
\Pr[\text{not } \varepsilon\text{-covered}] &\leq N(\varepsilon, \mathcal{M}) \cdot \left(1 - \frac{\varepsilon^{d_{\text{eff}}}}{D^{d_{\text{eff}}}}\right)^n \\
&\leq \left(\frac{3D}{\varepsilon}\right)^{d_{\text{eff}}} \cdot \exp\!\left(-n \cdot \frac{\varepsilon^{d_{\text{eff}}}}{D^{d_{\text{eff}}}}\right).
\end{align}

Setting this $\leq \delta$ and solving for $n$:
\begin{align}
n \geq \frac{D^{d_{\text{eff}}}}{\varepsilon^{d_{\text{eff}}}} \left[d_{\text{eff}} \ln\!\frac{3D}{\varepsilon} + \ln\frac{1}{\delta}\right].
\end{align}

For our setting with $D \approx 2$, $\varepsilon = 0.5$, $d_{\text{eff}} = 4$, $\delta = 0.05$:
\[
n \geq \frac{2^4}{0.5^4}\left[4 \ln 12 + \ln 20\right] = 256 \cdot [9.93 + 3.00] = 3{,}310.
\]

This bound is conservative; the empirical coverage at $n=50$ (coverage $\geq 0.87$) is much better because the data distribution concentrates on a low-dimensional manifold with favorable geometry. We use the empirical coverage with Clopper-Pearson confidence intervals as the primary certificate and the metric entropy bound as a theoretical complement.
\end{proof}

\subsection{Properties of Sinkhorn Divergence}

\begin{proposition}[Non-negativity and definiteness]
$\mathcal{S}_\varepsilon(\mu, \nu) \geq 0$ with equality iff $\mu = \nu$ (in the limit $\varepsilon \to 0$ for empirical measures).
\end{proposition}

\begin{proof}
This follows from the debiasing construction. The entropic OT cost satisfies $\text{OT}_\varepsilon(\mu, \nu) \geq \text{OT}_\varepsilon(\mu, \mu) + \text{OT}_\varepsilon(\nu, \nu) - \text{OT}_\varepsilon(\nu, \mu)$ doesn't directly give this; instead, $\mathcal{S}_\varepsilon$ is non-negative by construction as it equals the MMD with a Sinkhorn kernel, which is a positive definite kernel for $\varepsilon > 0$~\citep{feydy2019interpolating}.
\end{proof}

\begin{proposition}[Differentiability]
$\mathcal{S}_\varepsilon(\mu, \nu)$ is differentiable with respect to the positions of atoms in $\mu$ and $\nu$, with gradients given by the dual potentials.
\end{proposition}

This differentiability is what enables our greedy scoring: the gradient of the Sinkhorn divergence with respect to adding a new point is captured by the dual potential $g$ at that point's location.

\section{Extended Experimental Details}
\label{app:exp-details}

\subsection{Full Prompt List}

The 10 prompts used for real LLM evaluation:
\begin{enumerate}[leftmargin=*,itemsep=1pt]
\item \textbf{code\_review}: ``What are the most important things to check in a code review?''
\item \textbf{ml\_debug}: ``How do you debug a neural network that isn't learning?''
\item \textbf{system\_design}: ``Design a distributed cache for a social media platform.''
\item \textbf{math}: ``Explain three approaches to computing determinants of large sparse matrices.''
\item \textbf{ethics}: ``Key ethical considerations when deploying AI for hiring?''
\item \textbf{testing}: ``What's the best strategy for testing a complex microservice architecture?''
\item \textbf{security}: ``How do you perform a thorough security audit of a web application?''
\item \textbf{data\_pipeline}: ``Design a real-time data pipeline for processing IoT sensor data.''
\item \textbf{api\_design}: ``What principles matter most when designing a public REST API?''
\item \textbf{performance}: ``How do you diagnose and fix performance bottlenecks in a database-heavy application?''
\end{enumerate}

\subsection{Response Generation Parameters}

For each prompt, 8 responses are generated using GPT-4.1-nano with:
\begin{itemize}[leftmargin=*,itemsep=1pt]
\item 6 system prompts cycled: helpful, technical, creative, skeptical, practical, philosophical
\item Temperatures: 0.3, 0.5, 0.7, 0.9, 1.1, 1.3 (cycled)
\item Max tokens: 250
\item Total: $10 \times 8 = 80$ responses
\end{itemize}

\subsection{Embedding Statistics}

The \texttt{text-embedding-3-small} embeddings have dimension 1536.
Key statistics of the 80-response pool:
\begin{itemize}[leftmargin=*,itemsep=1pt]
\item Mean pairwise cosine similarity: $0.837 \pm 0.060$
\item Min cosine similarity: 0.656 (between responses to different prompts)
\item Max cosine similarity: 0.986 (between paraphrased responses to same prompt)
\item Effective dimension (95\% PCA variance): $\approx 4$
\end{itemize}

This confirms that LLM responses occupy a low-dimensional manifold within the 1536-dimensional embedding space, justifying our use of effective dimension in coverage certificates.

\subsection{Quality Score Computation}

Quality is computed as:
\begin{equation}
q_i = 0.4 \cdot \min\!\left(\frac{\text{len}(r_i)}{400}, 1\right) + 0.6 \cdot \min\!\left(\frac{\|x_i\|}{1.5}, 1\right),
\end{equation}
then linearly scaled to $[0.3, 1.0]$.
This is a simple proxy; in production systems, quality should use a reward model or human judgments.

\subsection{Kernel Bandwidth Selection}

For DPP and VCG-DivFlow, the RBF kernel bandwidth is set to the median pairwise Euclidean distance:
\[
\sigma = \text{median}(\{\|x_i - x_j\| : i < j\}).
\]
This is the standard median heuristic~\citep{kulesza2012determinantal}.
For the 80 LLM embeddings, $\sigma \approx 0.55$.

\subsection{Synthetic Experiment Setup}

Synthetic experiments use $N=200$ agents with Gaussian embeddings:
\begin{itemize}[leftmargin=*,itemsep=1pt]
\item Embedding dimension: $d = 16$
\item Agent types: Gaussian ($\mu \sim \mathcal{U}[-2, 2]^d$, $\Sigma = 0.5^2 I$), Clustered (4 clusters), Uniform, Mixture (2 components)
\item Quality: $q \sim \mathcal{N}(0.7, 0.1^2)$, clipped to $[0, 1]$
\item Seeds: 42, 137, 256, 512, 1024
\item Selection budget: $k = 10$
\end{itemize}

\subsection{Detailed Scaling Results}

\begin{table}[H]
\centering
\caption{Full scaling results: topic coverage and mean quality as pool diversity grows.}
\label{tab:scaling-full}
\small
\begin{tabular}{l|cc|cc|cc|cc|cc}
\toprule
& \multicolumn{2}{c|}{\textbf{DivFlow}} & \multicolumn{2}{c|}{\textbf{DPP}} & \multicolumn{2}{c|}{\textbf{MMR}} & \multicolumn{2}{c|}{\textbf{Random}} & \multicolumn{2}{c}{\textbf{Top-Q}} \\
\textbf{Pool} & Cov & Q & Cov & Q & Cov & Q & Cov & Q & Cov & Q \\
\midrule
2$\times$8 & 2/2 & .82 & 2/2 & .76 & 2/2 & .82 & 2/2 & .80 & 2/2 & .82 \\
4$\times$8 & 4/4 & .87 & 2/4 & .75 & 4/4 & .87 & 3/4 & .74 & 4/4 & .87 \\
6$\times$8 & 5/6 & .90 & 4/6 & .74 & 5/6 & .90 & 5/6 & .71 & 5/6 & .90 \\
8$\times$8 & 6/8 & .91 & 4/8 & .74 & 6/8 & .91 & 7/8 & .71 & 6/8 & .91 \\
10$\times$8 & 7/10 & .92 & 2/10 & .77 & 7/10 & .92 & 5/10 & .59 & 7/10 & .92 \\
\bottomrule
\end{tabular}
\end{table}

The scaling results (Table~\ref{tab:scaling-full}) reveal that DPP's coverage degrades dramatically at 10 topics (2/10, down from 100\% at 2 topics). This is because the DPP kernel becomes nearly diagonal in 1536 dimensions, causing greedy MAP to select high-quality items regardless of diversity.

DivFlow and MMR track closely in coverage and quality. The key differentiator is:
\begin{enumerate}
\item DivFlow provides the Sinkhorn dual potential interpretation (coverage certificates).
\item DivFlow integrates with VCG for incentive compatibility.
\item MMR's $\lambda$ parameter has a different interpretation than DivFlow's quality weight.
\end{enumerate}

\subsection{IC Verification Methodology}

For each IC trial:
\begin{enumerate}[leftmargin=*,itemsep=1pt]
\item Select a random agent $i$.
\item Compute truthful utility: $u_i^T = q_i \cdot \mathbf{1}[i \in S^*] - p_i$.
\item Generate a random deviation $\hat{q}_i \sim \mathcal{U}[0, 1]$.
\item Re-run the mechanism with $\hat{q}_i$; compute $u_i^D = q_i \cdot \mathbf{1}[i \in S'] - p_i'$.
\item Record violation if $u_i^D > u_i^T + 10^{-8}$.
\end{enumerate}

On real LLM data: 500 trials, 62 violations (12.4\%), max gain 0.90.
On synthetic data: 50 trials $\times$ 5 seeds, 9.2 mean violations (18.4\%).

\subsection{Kernel Comparison}

\begin{table}[H]
\centering
\caption{Kernel comparison on synthetic data (N=100, k=10, 5 seeds). Multi-scale kernel achieves highest diversity with lowest variance.}
\label{tab:kernels}
\begin{tabular}{lcc}
\toprule
\textbf{Kernel} & \textbf{Cosine Diversity} & \textbf{Log-Det Diversity} \\
\midrule
RBF ($\sigma=0.5$) & $0.989 \pm 0.033$ & $2.45 \pm 0.82$ \\
RBF ($\sigma=1.0$) & $0.992 \pm 0.028$ & $3.12 \pm 0.67$ \\
RBF ($\sigma=2.0$) & $0.985 \pm 0.041$ & $2.89 \pm 0.91$ \\
Adaptive RBF & $0.994 \pm 0.025$ & $3.28 \pm 0.58$ \\
Multi-scale & $\mathbf{0.997 \pm 0.019}$ & $\mathbf{3.41 \pm 0.44}$ \\
Manifold-adaptive & $0.993 \pm 0.027$ & $3.15 \pm 0.63$ \\
\bottomrule
\end{tabular}
\end{table}

\subsection{Scoring Rule Verification}

We verify properness of the scoring rules used in our mechanism:
\begin{itemize}[leftmargin=*,itemsep=1pt]
\item \textbf{Logarithmic rule}: $S(p, y) = \log p(y)$. Verified proper: 0 violations in 2500 tests.
\item \textbf{Brier rule}: $S(p, y) = 2p(y) - \|p\|^2$. Verified proper: 0 violations.
\item \textbf{Energy-augmented rule}: Adding a term $E(y, Y_{\text{hist}})$ that depends only on the outcome preserves properness. Verified: 0 violations.
\end{itemize}

\subsection{Pareto Frontier Analysis}

By varying the quality weight $\lambda$ from 0 (pure diversity) to 1 (pure quality), we trace the quality-diversity Pareto frontier:

\begin{table}[H]
\centering
\caption{Quality-diversity Pareto frontier (DivFlow, real LLM data).}
\label{tab:pareto}
\begin{tabular}{cccc}
\toprule
$\lambda$ & \textbf{Cosine Diversity} & \textbf{Mean Quality} & \textbf{Topics} \\
\midrule
0.0 & 0.987 & 0.716 & 1 \\
0.1 & 1.000 & 0.922 & 7 \\
0.5 & 1.000 & 0.922 & 7 \\
1.0 & 1.000 & 0.922 & 7 \\
\bottomrule
\end{tabular}
\end{table}

The sharp transition at $\lambda = 0.1$ indicates that a small quality weight suffices to move from pure diversity (selecting one extreme outlier) to balanced selection covering 7 topics.

\section{Replicator Dynamics: RPS Correction}
\label{app:rps}

The previous version of this paper reported that Rock-Paper-Scissors (RPS) converges to $(0.427, 0.299, 0.275)$ and marked it as ``Converged: Yes.''
This was incorrect.

\paragraph{Corrected analysis.}
The RPS payoff matrix $A = \begin{pmatrix} 0 & -1 & 1 \\ 1 & 0 & -1 \\ -1 & 1 & 0 \end{pmatrix}$ is a zero-sum game.
The interior equilibrium $(1/3, 1/3, 1/3)$ is a \emph{center}, not an attractor: the KL divergence $\sum_i x_i \log(x_i / (1/3))$ is a constant of motion under the continuous-time replicator equation.
Trajectories orbit around the equilibrium in closed curves.

With RK4 numerical integration ($dt = 0.01$), the orbits slowly expand due to numerical dissipation. After 1000 steps from $(0.4, 0.3, 0.3)$:
\begin{itemize}[leftmargin=*,itemsep=1pt]
\item \textbf{Converged}: No
\item \textbf{Equilibrium type}: Oscillatory center
\item \textbf{Mean of last 200 steps}: $(0.326, 0.284, 0.390)$ $\approx (1/3, 1/3, 1/3)$ with numerical drift
\item \textbf{Std of last 200 steps}: $(0.020, 0.003, 0.018)$ (oscillation amplitude)
\end{itemize}

\textbf{Corrected Table~\ref{tab:rps-corrected}:}

\begin{table}[H]
\centering
\caption{Corrected replicator dynamics results.}
\label{tab:rps-corrected}
\begin{tabular}{lccc}
\toprule
\textbf{Game} & \textbf{Converged} & \textbf{Final State} & \textbf{Theory} \\
\midrule
Prisoner's Dilemma & Yes & $(0.000, 1.000)$ & Defect dominates \\
Hawk--Dove & Yes & $(0.667, 0.333)$ & Mixed ESS $= (2/3, 1/3)$ \\
Rock--Paper--Scissors & \textbf{No} & \textbf{Oscillatory} & \textbf{Center at $(1/3, 1/3, 1/3)$} \\
\bottomrule
\end{tabular}
\end{table}

\section{Software Architecture}
\label{app:architecture}

The DivFlow framework consists of 60 Python modules organized into:

\begin{itemize}[leftmargin=*,itemsep=1pt]
\item \textbf{Core} (7): \texttt{mechanism.py}, \texttt{transport.py}, \texttt{kernels.py}, \texttt{dpp.py}, \texttt{coverage.py}, \texttt{scoring\_rules.py}, \linebreak \texttt{agents.py}
\item \textbf{Diversity metrics} (3): \texttt{diversity\_metrics.py}, \texttt{submodular.py}, \texttt{embedding.py}
\item \textbf{Mechanism design} (8): \texttt{mechanism\_core.py}, \texttt{auction.py}, \texttt{fair\_division.py}, \texttt{matching\_markets.py}, \texttt{voting\_systems.py}, \texttt{information\_elicitation.py}, \texttt{game\_theory.py}, etc.
\item \textbf{Extended modules} (10+): contract theory, network games, collective decisions, information design, etc.
\item \textbf{Tests}: 78 passing tests covering all core functionality
\item \textbf{Experiments}: Reproducible scripts with fixed seeds
\end{itemize}

\section{Reproducibility}
\label{app:reproducibility}

All experiments are reproducible:
\begin{itemize}[leftmargin=*,itemsep=1pt]
\item Synthetic experiments use fixed seeds: 42, 137, 256, 512, 1024
\item LLM experiments use deterministic API calls (fixed system prompts and temperatures)
\item Embedding model: \texttt{text-embedding-3-small} (deterministic)
\item All results are saved as JSON files in the repository
\item Test suite: \texttt{python3 -m pytest implementation/tests/ -q} (78/78 pass)
\end{itemize}

\end{document}
